\section{Agent interface and implementation}
\label{sec:implementation}
\label{sec:impl}

\subsection{MCP tools and lane policy}
\label{sec:impl-tools}

The platform exposes MCP tools across management and retrieval categories. They
fall into the management groups of Table~\ref{tab:functions}: parsing and
quality review for admission, library and sub-library creation, placement,
description and tag editing for organisation, move, rename, delete and forced
re-index for lifecycle, catalogue, vector, keyword and two-stage search for the
access path, the Laya evidence gate for segment admission, the document graph
for navigation, and operational-experience tools. These tools support a
closed operational loop: a write changes managed state, a read or probe checks
the resulting state, and a re-index or reconciliation operation makes stale
state visible or repairable. Persona-distillation workflows are outside the
scope of this evaluation and use their own retrieval and synthesis path.
Access policies receive read-oriented tool sets; administrative writes remain
separate, and long-running parse, re-index and graph operations return
pollable task identifiers. Figure~S1 summarises the software design: a tool
layer over management services that share one auditable state, with the
evidence gate as an optional segment-level call.



The chat endpoint adapts multiple agent harnesses to the lane interface. The
server selects a harness and permission mode, enforces configured tool
allow/deny lists, applies endpoint timeouts and turn limits, and may register an
asynchronous task. Lane-specific retrieval sequences and call budgets are
workflow instructions; the endpoint does not enforce a separate per-lane call
quota.

\subsection{Runtime and deployment}
\label{sec:impl-runtime}

The backend uses Python 3.12 and FastAPI. ChromaDB stores cosine-space HNSW
collections with normalised BGE-M3 embeddings (1024 dimensions)
\cite{chroma,hnsw,bge-m3}; Neo4j stores the document graph \cite{neo4j};
BM25 supports lexical search \cite{bm25}; MinerU provides parsing and OCR
\cite{mineru}. The Laya verifier runs locally as a serialised worker.

The evaluation environment is one self-hosted Windows workstation with an
RTX 5060 (8 GB) used for embeddings and judging. Local execution reduces
external service dependencies but does not establish privacy, access-control or
multi-host deployment guarantees. A stack-ensure script checks the API, web
front end and MCP endpoints and can restart missing services before an
end-to-end retrieval probe. Three recent acceptance trials reported recovery
from a stopped stack without manual intervention. This is an operational
observation, not a guarantee of availability in every deployment.

The implementation and scripts are available
at \url{https://github.com/kingdol666/rag-knowledge}; experimental corpora
remain subject to their data-use constraints.
